\chapter{DISLOCAMENTO}\label{dislocamento}

\hfill\break

Alla fine della primavera successiva ero in Ecuador, sul punto di essere
spedito oltre i confini del mondo per combattere i criceti nello spazio
profondo. Non è un calcio in culo? Ero sorpreso di quanto facesse fresco
sulle montagne dell'Ecuador, le mie immagini mentali del Sud America
avevano sempre a che fare con il caldo soffocante, come l'avevo
sperimentato in Nigeria. Era stato un inverno freddo in Maine, non
freddo in termini di temperature e neve, che erano state nella media
degli ultimi anni. Freddo in termini di niente elettricità per un lungo
periodo di tempo e scarse forniture di benzina e gasolio da
riscaldamento. La mia città natale e la mia gente erano meglio
attrezzati per l'inverno di molti altri in tutto il paese: i miei
genitori avevano una stufa a legna in casa e un'altra nel
garage-officina. Una famiglia della regione meridionale, che era stata
cacciata dal suo condominio, venne a vivere nel garage dopo Natale, mio
padre barattò la riparazione del motore di un trattore con delle balle
di paglia che sistemammo sulle pareti del garage come isolante,
rivestendole di tela cerata. Feci visita ai miei genitori un paio di
volte, il garage era accogliente e caldo e odorava di fieno appena
falciato. C'era un'altra famiglia di tre persone, più una donna single,
che viveva nella casa dei miei genitori, nell'ambito di un programma di
reinsediamento governativo.

Tutti fecero del loro meglio, per tutto l'inverno.

I governi statali e federale degli Stati Uniti ebbero un inizio stentato
nel fronteggiare la crisi dopo l'attacco ruhar, mettendo tutto
sottosopra invano, poi si coordinarono e si concentrarono solo su come
far sì che le persone passassero indenni l'inverno. L'attacco ruhar non
era stato quello che ci saremmo aspettati dalla visione di un film di
fantascienza: invece di colpire le basi militari e i principali centri
abitati, avevano per lo più distrutto centrali elettriche e impianti
industriali. Era strano: vidi le foto satellitari di New York e
Washington dopo l'attacco e, a parte la mancanza di luci e traffico per
le strade, le città apparivano intatte. Dopo che i kristang avevano
cacciato i ruhar dal nostro sistema solare, gli Stati Uniti d'America
non avevano molte infrastrutture elettriche funzionanti. Internet e reti
telefoniche erano inattivi, la fornitura di elettricità era irregolare,
sempre che ce ne fosse, la radio trasmetteva soltanto messaggi di
emergenza un paio di volte al giorno, stazioni Tv e telegrafo erano
fuori uso. La Marina collocò portaerei e sottomarini nucleari nei porti
delle principali città costiere e li collegò in funzione di centrali
elettriche galleggianti, per fornire elettricità a strutture critiche
come gli ospedali. Stavamo lentamente riguadagnando energia elettrica in
tutto il paese, sottolineo ``lentamente''. Quando partii per l'Ecuador,
i miei genitori non avevano ancora energia, a parte quella prodotta dal
generatore che mio padre aveva collegato alla presa di potenza del
trattore. Faceva funzionare il mezzo con una miscela in parti uguali di
alcol preparato in casa e benzina. Il primo consumava le guarnizioni del
motore, perciò mio padre e il tizio che viveva in garage le sostituivano
ogni mese. Mettevano in funzione l'elettricità solo la mattina presto e
la sera, comunque, non c'era carburante disponibile per molto di più.

Il problema più grande che l'America stava affrontando non era la
mancanza di elettricità, ma la crisi economica seguita all'attacco. Mio
padre perse il lavoro quando la cartiera chiuse. I boscaioli non avevano
abbastanza benzina per far funzionare i loro camion e anche le loro
motoseghe. Senza la pasta di legno, la cartiera non poteva produrre
carta. Questo mercato era comunque crollato insieme al resto
dell'economia globale. Nazioni industriali come gli Stati Uniti, il
Giappone, la Cina, e tutta l'Europa furono colpite in modo molto più
duro delle aree a basso sviluppo tecnologico del pianeta, ironia della
sorte. Maggiore era tale livello in una nazione, più forte questa era
colpita dalla crisi. Il valore delle aziende tecnologiche finì dritto
nel cesso, non solo a causa della mancanza di domanda da parte
dell'economia generale. Chi voleva investire nella Silicon Valley,
quando i nostri nuovi alleati, i kristang, avrebbero condiviso con noi
tecnologie incredibilmente avanzate?

Non lo fecero però. I kristang dissero che non eravamo pronti, che non
potevano affidarci il loro livello tecnologico e che dovevamo
concentrarci sulla ricostruzione della nostra infrastruttura attuale.
Davvero, i nostri nuovi alleati furono una delusione, a parte il fatto
di avere cacciato i ruhar. La maggior parte del gruppo di battaglia
kristang se ne andò in fretta e furia nel giro di una settimana dopo
aver sconfitto i ruhar, perché la Terra non aveva porti spaziali o
strutture di servizio per le navi da guerra, o qualsiasi altra cosa di
cui i kristang avessero bisogno. Non combattevano i ruhar per il bene
dell'umanità, combattevano per negare ai ruhar una base nel nostro
piccolo angolo di galassia. La Terra era la Guadalcanal della seconda
guerra mondiale, così la descrisse uno degli ufficiali della guardia
nazionale. Nessuna delle due parti si preoccupava del luogo, se non come
trampolino per un posto più importante. Gli Stati Uniti e il Giappone
non si erano preoccupati dell'isola di Guadalcanal o dei nativi che
vivevano lì, l'unica cosa che entrambe le parti volevano era usare
quell'isola come base per spingersi verso la successiva conquista. Così
era la Terra per i kristang e i ruhar. Essendo figli dell'America,
l'unica superpotenza al mondo, la più ingente forza militare della
storia, era difficile per noi immaginarci nel ruolo dei primitivi
nativi. Ce ne stavamo a guardare le forze kristang e ruhar combattere
sulla nostra terra, con armi che riuscivamo a malapena a capire.

\hfill\break

«Ehi! Bishop! Ehi!»

Girai su me stesso, disorientato, per ritrovarmi faccia a faccia con un
tipo del mio gruppo di fuoco, un ragazzo che non avrei mai pensato di
rivedere. «Cornpone!», gridai.

«Brown Bread!», rispose, ci abbracciammo e ci prendemmo a pacche sulla
schiena con una tale forza che restai senza fiato. Jesse Colter,
dall'Arkansas e orgoglioso figlio del Sud, come mi aveva detto quando
l'avevo conosciuto durante l'addestramento di base. L'avevo chiamato
Cornpone, cioè ``pane di mais'', perché era, credo, qualcosa che la
gente mangiava al Sud. Mi aveva risposto per le rime chiamandomi ``pane
nero'', in riferimento a un pane in lattina, addolcito e farcito di uva
passa, che è una tradizione del New England. Una volta gliene avevo
servita una fetta. Alcune persone mettono in acqua calda la lattina, ma
io preferisco tostare ogni fetta. Si tagliano entrambe le estremità
della lattina, che incide anelli sul pane mentre lo si spinge fuori.

Lo so. Fidatevi, è delizioso.

«O adesso dovrei chiamarti Barney?», rise Jesse.

Merda. La storia si era diffusa, i ragazzi della guardia nazionale con
cui ero mi avevano sfottuto per tutto l'inverno. ``Barney e i Puffi'', è
così che la gente chiamava me e la banda che aveva catturato il soldato
alieno. «Oh, amico, lo sanno tutti adesso?»

«Col supporto di internet, l'intero pianeta lo sa. Accidenti, amico,
bello vederti!», disse Jesse. «Ce ne sono un po' del X qui, ma ancora
nessun altro della nostra squadra.»

«Sei la prima persona del X che vedo», ammisi, guardandomi attorno alla
ricerca di altre facce familiari.

«Come sei arrivato qui?», chiese Jesse, chinandosi per prendere il
borsone che stava trasportando prima di vedermi.

«Dopo il Giorno di Colombo, stavo cercando di tornare a Drum, ma nessuno
aveva abbastanza benzina per fare il viaggio. Perciò, mi sono messo in
contatto con la guardia nazionale locale, chiedendo se sapevano cosa
stesse accadendo. Il colonnello della guardia mi ha obbligato
all'azione, erano stati federalizzati comunque. Mi sono detto: ``Che
cazzo, mi hanno dato un fucile, un casco e del cibo'', giusto?» Gli
ordini di assegnarmi ufficialmente alla guardia erano arrivati con la
radio a onde corte più tardi. «Così, ho fatto il soldato nel Maine, ho
aiutato nella raccolta, nella sorveglianza dei convogli di carburante e
cibo. Poi ho ricevuto la chiamata per venire qui, così sono saltato su
un treno per Boston.» Un treno merci. Viaggiando su un vecchio carro
merci convertito, mi ero sentito come un vagabondo, anche se avevo un
biglietto del governo degli Stati Uniti.

«No, intendevo, come siete arrivati qui? In volo?»

«Oh, sì, ehm, ci hanno messo su una tradotta da Boston, ci sono voluti
cinque giorni per arrivare a Miami. Siamo atterrati circa due ore fa.»
Avevamo volato su uno United 767 che aveva visto giorni migliori.

Jesse fece di sì con la testa. «Ho preso un treno per Dallas e me ne
sono rimasto lì seduto tre giorni ad aspettare un volo, la cazzo di
Aeronautica non riusciva a racimolare abbastanza carburante per fare
benzina sull'aereo. Ho sentito che alcune unità stanno arrivando sulle
navi da New Orleans e Houston.»

«Hai visto questa specie di torre?»

«L'ascensore? Solo quello che si può vedere da qui. Non ne so niente,
compare.»

Alzammo gli occhi alla montagna, la cima era appena visibile sotto un
cielo coperto. Sparare direttamente da lassù era un gioco da ragazzi.
Quando l'Onu aveva accettato di fornire truppe per combattere sotto il
comando dei kristang, i kristang avevano costruito quello che definivano
un ascensore spaziale, realizzato in meno di un mese. Avevano scelto
quella montagna in Ecuador perché si trovava a cavallo dell'equatore.
Sulla cima, i kristang avevano costruito una stazione base con un
reattore a fusione e, migliaia di chilometri sopra, c'era una stazione
in orbita geosincrona. Ciò che collegava le due stazioni era un sottile
campo magnetico, in sostanza un fulmine stazionario. Dovevamo entrare in
orbita con un ascensore che percorreva il fulmine stazionario. Uno dei
ragazzi dell'Aeronautica sul mio volo disse di aver sentito che i
kristang inviavano nel campo magnetico degli impulsi che spingevano
l'abitacolo dell'ascensore verso l'alto. Quell'ascensore in Ecuador era
stato il primo a essere terminato, i kristang parlavano di costruirne
altri due, uno in Africa e l'altro da qualche parte in Indonesia.
Strizzai gli occhi e mi feci scudo con una mano, esaminando il sottile
fascio di luce. «Non so neanche di questo, amico. Una cosa che so è che,
se i kristang sono in grado di fare una cosa del genere, sono felice che
siano dalla nostra parte.»

«Ti capisco, amigo», concordò Jesse.

Guardando il cielo, il fulmine che avrei percorso per andare in orbita,
mi chiesi che ne fosse stato di Faccia Pelosa. Ci eravamo diretti a sud
col nostro furgone dei gelati, finché non ci eravamo imbattuti in un
posto di blocco della polizia e, quando avevano visto chi c'era nel
retro del furgone, ci avevano fornito una scorta fino a Lincoln, dove un
elicottero della guardia nazionale aveva portato via Faccia Pelosa. Dopo
di che, ogni informazione su di lui era riservata ed ero stato dissuaso
dal parlare o chiedere di quel criceto. Ovunque fosse, speravo che lo
trattassero bene.

Un suono ronzante catturò la nostra attenzione, ci girammo entrambi per
guardare un gruppo di convertiplani grigi che volavano in formazione e
provenivano da ovest. Mentre si avvicinavano, gli aerei di testa
rallentarono, ruotarono i rotori nel passaggio al volo a punto fisso,
poi calarono come elicotteri. Prima che potessi chiederlo, Jesse disse
spontaneamente: «La Marina ha due gruppi di portaerei al largo della
costa, il Lincoln e il Reagan. Ho incontrato alcuni Marine qui nei
dintorni stamattina».

«Vengono con noi?» Due gruppi di portaerei sarebbero sembrati imponenti,
prima che i ruhar attaccassero. Rispetto alle forze armate aliene, che
potevano sollevare le persone nello spazio con un fulmine, le nostre
portaerei a propulsione nucleare erano però barche a remi obsolete.

«Ci sono alcune unità del corpo dei Marine, ma questi tipi? No, sono qui
per mantenere l'ascensore in sicurezza.»

Avevo sentito alla radio che i kristang avevano affidato la sicurezza
dell'ascensore spaziale in Ecuador agli Stati Uniti. Risi. «Scommetto
che ne sono felici.» Nessun Marine americano che si rispetti sarebbe
voluto restare di guardia mentre l'esercito andava nello spazio per
vendicare l'attacco al nostro pianeta.

«Cazzo sì, amico. Mi hanno offerto un sacco di soldi per questo
distintivo.» Jesse si tastò lo stemma blu dell'Unef sulla manica destra,
il logo della Expeditionary Force delle Nazioni unite, come ci chiamano
ora. Gli stemmi nazionali, nel nostro caso la bandiera degli Stati
Uniti, andavano sulla manica sinistra. «Come se il denaro significasse
tutto ora.» Jesse si accigliò e abbassò gli occhi sui propri abiti.
«Almeno quei Marine hanno delle vere uniformi. Questa l'ho presa in un
negozio in Arkansas.» I pantaloni gli stavano giusti, mentre la giacca
sembrava di tre taglie troppo grande.

Anche la mia ``uniforme'' era spaiata: i pantaloni erano della guardia
nazionale del Maine, mentre la parte di sopra, che mi era stata data sul
treno per Miami, era di un vecchio modello mimetico e sembrava fosse
rimasta in deposito dalla guerra in Iraq, la prima. O forse dalla guerra
ispano-americana. Emanava ancora un fievole odore di naftalina. «Forse
l'esercito ci assegnerà nuove uniformi qui.»

«Non ci contare.»

Ci abbaiò contro una voce familiare. Ci girammo e vedemmo il tenente
Amos Gonzalez, un capo plotone nel X. Non il nostro plotone, ma un volto
conosciuto. Jesse e io scattammo sull'attenti e salutammo. «Riposo,
uomini. Bishop e Colter, giusto?» La sua faccia si aprì in un ampio
sorriso. «È bello vedervi entrambi. Contento che ce l'abbiate fatta ad
arrivare alla festa.»

«Anch'io sono felice di vederla, tenente», disse Jesse. «Sa quanti di
noi sono qui?»

L'uomo scosse la testa abbassando gli occhi, mentre calciava la terra
con lo stivale. «Le comunicazioni sono ancora confuse. Stiamo mettendo
le persone al lavoro a mano a mano che arrivano alla spicciolata.»

Jesse si accigliò guardando Gonzalez, che aveva aggirato la domanda.
«Ehi, tenente, finora ho visto solo noi, i ragazzi dei Belching
Buzzards¹», intendeva la CI Screaming Eagle², «il III fanteria e i
Marine. Non portiamo corazzati o cavalleria?»

Il tenente Gonzalez gli gettò un'occhiata sprezzante. «Colter, davvero
vorresti stare dentro un carro armato a combattere contro un nemico che
può sparare dall'orbita?» Non aspettò la risposta di Cornpone. «È un
passaporto per l'inferno. I kristang hanno detto che avevano bisogno di
fanteria, quindi questo è quello che stiamo inviando. Corazzieri e
aviatori resteranno in panchina a questo giro.»

«Mi perdoni, signore, ma a che diavolo serve la fanteria contro i
ruhar?», chiesi. «Non ho nemmeno un'arma.»

Gonzalez annuì con simpatia. «I kristang non permetteranno che si
portino armi sull'ascensore. Dovremo procurarci armi, giubbotti
antiproiettile e tutto il resto quando arriveremo, ovunque andremo. Per
quanto mi riguarda, mi piace il mio vecchio M4, ma non voglio andare con
questo contro i ruhar.»

«Siamo in servizio coi Marine, signore?», chiesi.

«Una brigata del corpo dei Marine, sì. E, ho sentito, inglesi, francesi,
cinesi, forse anche alcuni russi e l'esercito indiano. Avete già fatto
il controllo medico?»

«No, signore, sono appena sceso dall'autobus proveniente
dall'aeroporto», dissi. Riuscivo ancora a sentire le doghe d'acciaio del
sedile dell'autobus sotto il mio sedere. Eravamo arrivati
dall'aeroporto, ci eravamo inerpicati per la strada sterrata, a stento
segnata, in un convoglio di vecchi scuolabus dai colori sgargianti: è
probabile che fossero i migliori mezzi di trasporto che quella parte
dell'Ecuador aveva da offrire prima che i ruhar attaccassero.

«Ehm, ehm.» Jesse scosse la testa. «Sono qui da stamattina e, appena
sono sceso dall'autobus, i Marine mi hanno fatto spostare bancali di
cibo.»

«Fareste meglio a fare un controllo, allora, avete bisogno di una
valutazione medica prima d'imbarcarvi sul tappeto magico per lo spazio.
Vedete quel grande tendone da circo laggiù, quello con le bandiere delle
Nazioni unite?» Non scherzava, era davvero un tendone da circo.
«Andateci di corsa, l'ascensore parte alle 17:00.»

\hfill\break

La tenda era un manicomio. Jesse e io non eravamo impazienti di farci
punzecchiare e pungolare dai medici dell'esercito, ma non avremmo avuto
ragione di preoccuparci. L'esame medico consisteva nello stare in piedi
vestiti, in una cabina delle dimensioni di una doccia, mentre facevano
scendere e salire e di nuovo scendere un fascio di luce. Quando uscii,
il dottore non alzò gli occhi dal suo computer. «Vivrai, Bishop.» Feci
in tempo a vedere un file con la mia immagine sullo schermo del
computer, il medico pigiò alcuni tasti e il file fu sostituito da quello
di un altro ragazzo. «Avanti il prossimo!»

«Mi scusi, signore, è tutto qui?», chiesi confuso.

«È tutto qui. Hai appena fatto il tuo primo esame con uno scanner medico
kristang e dice che sei abbastanza in salute. Su, gente, non abbiamo
tutto il giorno. Avanti il prossimo!»

Dopo l'esame medico, facemmo la doccia, con autentica acqua calda, anche
se non ci assegnarono nuove uniformi, così ci rimettemmo i vecchi
vestiti e poi andammo in un'altra tenda per un pasto caldo, stile
esercito. Per quanto mi riguardava, lo stile militare era una buona
cosa, significava che avrei avuto un sacco di cibo. Attraversai la fila,
ero in piedi con il mio vassoio in mano in cerca di un posto per
sedermi, quando Jesse si alzò e mi fece cenno di andare al suo tavolo.
«Questo qui è Gus, I corazzato», disse Jesse con la bocca piena di pane
di mais, indicando un ragazzo seduto accanto a lui.

«Ma non doveva esserci solo la fanteria in questo viaggio?», chiesi.

Gus alzò le spalle. «Ho guidato un Bradley, immagino che abbiano bisogno
di autisti per qualsiasi cosa ci sia lassù.» Se i kristang pensavano che
ci servissero autisti, si poteva sperare che noi della fanteria non
saremmo stati costretti a farcela sempre a piedi ovunque fossimo andati.

Cornpone usò un pezzo di pane per tirare su il sugo dal polpettone. «Non
so dove stiamo andando, ma, se significa tre bocconcini e una branda, ci
sto. Era piuttosto magra quest'inverno a casa.»

\hfill\break

Non sapevo cosa aspettarmi dall'ascensore spaziale, ma non quello che mi
si presentò davanti. Sembrava più la sala d'attesa di un aeroporto.
L'enorme abitacolo dell'ascensore era a forma di disco e aveva quattro
livelli, il più basso dei quali conteneva i macchinari. O così supposi,
perché ci fu consentito l'accesso soltanto ai tre livelli superiori.
Capsule di carico erano agganciate sotto la parte inferiore
dell'abitacolo. I livelli abitabili erano divisi ciascuno in otto
settori, contenenti file e file di sedie di plastica dall'imbottitura
sottile. Sedie con cinture di sicurezza. La durata prevista per il
viaggio fino alla stazione orbitale era di nove ore, ore movimentate
solo dal nostro affollarci attorno agli oblò per guardare giù, verso la
superficie del nostro pianeta che si allontanava. Dopo un po', anche
quello ci stufò. L'ascensore era un aggeggio senza fronzoli, progettato
per trasportare un numero massimo di truppe in un periodo minimo di
tempo. Il comfort durante il viaggio non era una priorità dei kristang.
Non c'era nemmeno uno snack bar integrato nell'ascensore, quindi
l'esercito aveva allestito una cucina mobile a ogni livello e il cibo
era abbondante. Mangiai il mio primo cheeseburger decente da molto tempo
a quella parte, poi tornai indietro a prenderne un altro. Jesse e io
esplorammo il posto, girovagammo, incontrammo i soldati con cui avremmo
prestato servizio, condividemmo storie, ascoltando voci contraddittorie.
Poiché l'ascensore aveva una capacità di cinquemila umani, buona parte
di una divisione dell'esercito era a bordo durante il nostro viaggio.
Passammo accanto a un settore che era stato adattato a mensa provvisoria
degli ufficiali, ci infilammo la testa, per poi proseguire e mescolarci
con gli altri soldati semplici.

Dopo essermi scolato tre soda -- l'esercito non forniva bevande
alcoliche -- ebbi bisogno di fare una capatina in bagno, cosa che
temevo. Che tipo di tubature avevano fornito i kristang? Beh, o questi
ultimi avevano funzioni biologiche simili alle nostre, o avevano
consultato idraulici umani, perché in realtà la mia preoccupazione si
rivelò immotivata. Sembrava un tipico bagno da autogrill, solo nuovo e
più pulito. C'era anche il sapone nei dispenser.

Il viaggio all'inizio fu lento, era davvero come andare in ascensore,
c'erano una sottile pressione provocata dalla gravità in eccesso e una
lieve vibrazione. Mentre salivamo e l'atmosfera si rarefaceva,
accelerammo e la gravità artificiale si attivò per compensare, a quello
che ci dissero essere circa l'85\% del normale sulla Terra. Parlai con
uno dei cuochi dell'esercito, che aveva fatto su e giù con l'ascensore
diverse volte. Mentre ci avvicinavamo alla stazione spaziale, spiegò,
rallentavamo, ma non sentivamo alcun movimento a causa della gravità
artificiale. Espressi dubbi sulla sicurezza di percorrere un fulmine,
soprattutto con i predoni ruhar che ancora attaccavano i kristang, il
cuoco mi disse di non preoccuparmi. All'inizio della settimana, una
fregata ruhar era saltata fuori dall'iperspazio mentre l'ascensore era a
metà strada per la stazione spaziale, nel suo punto più vulnerabile. La
fregata era riuscita a lanciare due missili prima che un paio di
cacciatorpediniere kristang la colpissero e che laser difensivi
dell'ascensore si occupassero dei missili in arrivo. Fu bello sapere che
non eravamo proprio bersagli facili, mentre mangiavamo spuntini e
stavamo seduti su sedie scomode.

Riuscii a addormentarmi, spaparanzato su una sedia. Mi svegliai quando
un Marine mi diede un calcio al piede. «Alzati e splendi, bella
addormentata. Questo giro in auto rubata sta arrivando alla fine.
Raccogli la tua attrezzatura, se ne hai una.» Il Marine si allontanò per
andare a svegliare a calci altri soldati addormentati.

Jesse si era assopito sul pavimento dietro di me, usando il suo borsone
come cuscino. Gli strinsi la spalla. «Oh, amico, la tua brutta faccia
non è quello che voglio vedere svegliandomi», esordì scontroso. «Come
va?», riuscì a dire, mentre uno sbadiglio gli stirava la mascella.

«Credo che ci stiamo avvicinando alla stazione spaziale. Vuoi dare
un'occhiata?»

\hfill\break

La stazione era una ciambella, con sei braccia per raggi. All'estremità
di due delle braccia c'erano quelle che diedi per scontato fossero navi
kristang, a riposo a pancia in giù, così sembrava, su piattaforme. Una
delle astronavi era elegante e scura, l'altra molto più grande e non
così potente. Mentre io e Jesse le fissavamo inebetiti, il cuoco con cui
avevo parlato prima si fece largo tra la folla per avvicinarsi a noi.
«Quello», indicò la nave più piccola, «è un incrociatore, e quella
brutta e grossa accanto è la nave da trasporto.»

«Ci sei salito?», chiese Jesse.

«No, ci fanno salire sulla stazione solo mentre l'abitacolo
dell'ascensore sta imbarcando, poi si torna in Ecuador per il gruppo
successivo. Ho sentito che la nave da trasporto può contenere l'intero
gruppo, non appena sarete a bordo partiranno.»

«Hai idea di dove ci stanno mandando?» Pensavo che il cuoco potesse
avere sentito qualcosa durante uno dei suoi viaggi precedenti.

Lui scosse la testa. «Non lo so, e se qualcuno dei nostri ufficiali lo
sa, non lo dice. Per quanto mi riguarda, penso che i kristang informino
soltanto gli addetti ai lavori, e noi umani non lo siamo.» Guardammo in
silenzio, mentre l'ascensore si avvicinava lentamente alla stazione,
finché la massa di quest'ultima gettò un'ombra attraverso il finestrino
e nascose l'incrociatore alla vista. «Vorrei tanto venire con te,
provare a passare alla fanteria», il distintivo sulla sua uniforme
diceva I corazzato, «ma l'esercito dice che il mio ruolo è quello di
cuoco, quindi questo è quello che sto facendo.» Mi sembrò amareggiato.

«Ehi, il tuo cheeseburger è la fine del mondo.» Gli tesi la mano.

Lui si strofinò la sua sul grembiule, afferrò la mia con forza e mi
guardò dritto negli occhi. «Voi ragazzi state per spaccargli il culo,
vero?»

``Puoi giurarci'', risposi con convinzione. Non sapevo se sarei tornato,
ma avrei ucciso il maggior numero possibile di criceti. Ogni soldato e
Marine con cui parlavo provava la stessa cosa, lo stesso mortale senso
di determinazione. Eravamo stati attaccati, la sopravvivenza stessa
della nostra specie era stata minacciata. ``Rabbia'' non era la parola
adeguata a descrivere come ci sentivamo tutti. Nella stazione
ferroviaria di Boston camminavo con altri tre soldati diretti verso lo
spazio e la guerra, e la folla ci aveva notato. La gente si fermava,
iniziava spontaneamente a salutare, poi applaudiva. Ero davvero
commosso. Ricordo che un vecchio, doveva essere un nonno, aveva con sé
il nipotino. Aveva mostrato al piccoletto come stare sull'attenti e gli
aveva sistemato la mano in un saluto militare rivolto a noi.
L'espressione sul viso del nonno, proprio in quel momento, mi aveva
detto tutto di quello che la gente provava per noi soldati della
Expeditionary Force delle Nazioni unite. Stavamo portando con noi tutte
le loro speranze, non solo le speranze di vendetta. Speranza per la
sopravvivenza, la sopravvivenza della specie umana. È ciò che rendeva
diversa quella guerra: per la prima volta eravamo davvero tutti dalla
stessa parte.

«Puoi giurarci, amico, stiamo per spaccargli il culo», disse Jesse in
tutta serietà, e pigiammo pugno contro pugno. «Esercito forte non va
alla morte. Urrà!»

\hfill\break

Come il cuoco, non vedemmo molto dopo avere lasciato l'abitacolo
dell'ascensore. Le guardie del corpo dei Marine ci mandarono alla
stazione spaziale che generava il fulmine da noi percorso in direzione
dello spazio. Ci mettemmo in fila indiana, per fare cosa non lo sapevo.
La stazione era grande, mi chiedevo ad alta voce come avessero fatto a
costruirla così in fretta.

«Una nave madre dei thuranin l'ha portata qui in sezioni, poi l'hanno
assemblata», disse un tenente in fila davanti a noi.

«Signore?», chiesi. «Cosa sono i thuranin?»

«I thuranin sono la specie patrona dei kristang.» Reagendo agli sguardi
del tutto vuoti che gli rivolgemmo io, Cornpone e tutti gli altri, uscì
un po' dai ranghi per parlare con noi. Dalla sua uniforme, vidi che era
del III fanteria. «Non vi hanno detto niente?»

«Signore, vengo dal Nord del Maine, abbiamo avuto internet per la prima
volta una settimana prima dell'attacco dei ruhar.» Vidi che non
apprezzava il mio tentativo di fare dell'umorismo. «Non ho mai sentito
parlare dei thuranin, signore.»

Lui fece roteare il dito in aria. «Questa stazione, l'ascensore, è tutta
tecnologia thuranin, i kristang non hanno nemmeno la gravità artificiale
a bordo delle loro navi.»

«Signore?» La faccia di Jesse iniziò a diventare verde quando lo sentì.
Tutti nella fila apparivano nauseati all'idea di sperimentare l'assenza
di gravità. «Gravità zero? Non siamo stati addestrati per questo, per
niente.»

«Ecco perché sei in fila qui, devi prendere le medicine per prevenire la
nausea spaziale. Non possiamo permettere che gli umani vomitino le
budella su una nave da trasporto kristang.»

«Non andremo con un'astronave dei thuranin?», chiesi, sentendo già la
nausea. Mi stavo pentendo amaramente di quel secondo cheeseburger. La
fila avanzava e noi strascicavamo i piedi lungo il ponte.

«Le navi madri thuranin si avventurano di rado in un pozzo
gravitazionale, rimangono ben lontane nella Nube di Oort di qualsiasi
sistema stellare.» Vedendo altri sguardi vuoti, sebbene io mi
compiacessi di sapere che cosa fosse una Nube di Oort, aggiunse: «Questo
significa fuori dall'orbita di Plutone, dove lo spazio-tempo è piatto ed
è più facile per le loro navi madri formare un campo di salto. Le navi
kristang possono saltare solo a corto raggio, tra le sei e le otto ore
luce, che è più o meno fino a Nettuno. Fanno una serie di piccoli salti
più veloci della luce in questo modo, fino a dove è parcheggiata una
nave madre thuranin. Una nave madre è una grande, lunga colonna
vertebrale, con punti d'attacco per agganciare le astronavi a breve
raggio. Ha un motore di salto avanzato, ben oltre quello dei kristang,
quindi i thuranin trasportano le navi kristang tra i sistemi stellari.
La G2, la nostra intelligence a livello di divisione, pensa che le navi
kristang non possano di fatto viaggiare tra le stelle: con i loro salti
brevi e la necessità di ricaricare i motori tra un salto e l'altro, nel
complesso viaggiano a velocità minore di quella della luce. Ci
vorrebbero più di quattro anni per arrivare da qui alla stella più
vicina. I thuranin, pensiamo, possono viaggiare a qualcosa come trecento
volte la velocità della luce, in salti di circa un anno luce».

Non mi piacevano i calcoli matematici che mi stavano frullando nel
cervello già quasi fumante. Trecento volte la velocità della luce era
incredibilmente veloce, ma significava anche che, se fossimo andati su
un pianeta che era, diciamo, a un solo centinaio di anni luce di
distanza, ci sarebbero voluti quattro mesi per arrivarci! Quattro
maledetti mesi. Bloccato su una nave aliena. In assenza di gravità. Un
centinaio di anni luce sembrava un sacco, ma la galassia è enorme. A un
certo punto a scuola avevo imparato che la Terra si trova a
ventisettemila anni luce dal centro della galassia, e l'intera galassia
è qualcosa come centomila anni luce di diametro. Non avevo considerato
nulla di tutto questo, ansioso com'ero di lasciare la Terra per
combattere i ruhar. «Signore», chiesi lentamente, «dove stiamo andando?»

«Non è un segreto. Stiamo andando in un pianeta che i kristang hanno
adibito a base di addestramento per noi, lo chiamiamo Campo Alfa.»

«E quanto dista?», chiese Cornpone gettandomi uno sguardo. Stava
calcolando mentalmente anche lui.

«I kristang non ce lo diranno, ma le persone che ci sono state hanno
scattato foto del cielo notturno e, dalla posizione delle stelle, i
nostri astronomi l'hanno fissato a 1.223 anni luce dalla Terra.» Proprio
allora, una delle porte di fronte alla fila si aprì e il tenente fu
invitato a entrare.

Oh, mio Dio. Un viaggio di milleduecento anni luce, a bordo di un mezzo
che viaggiava a trecento volte la velocità della luce, significava che
saremmo rimasti bloccati a bordo di una nave kristang, agganciata a una
nave madre thuranin, in assenza di gravità, per quattro anni! Avevo
sentito dei problemi che gli astronauti avevano avuto a bordo della
Stazione spaziale internazionale a gravità zero, di come la loro vista
si fosse deteriorata, le loro ossa si fossero indebolite e i loro
muscoli fossero deperiti dopo un paio di mesi. Che cosa avrebbero potuto
mai fare di buono dei soldati umani dopo quattro anni a gravità zero?
Ero così scioccato, che non mi venne neppure in mente che il tenente
aveva detto che in qualche modo gli umani erano già andati al Campo Alfa
e ne erano tornati, ovviamente in meno di un anno.

«Quattro anni, cazzo! Giusto? Milleduecento anni luce di distanza, a
trecento all'anno, no?» La voce di Cornpone rifletteva il turbamento che
tutti noi avvertivamo. «Forse ci congeleranno, ci faranno dormire per
tutto il tragitto? L'hanno fatto in Avatar, giusto\ldots{} hai visto
quel film?»

La gente mormorava, stordita dalle implicazioni di ciò che il tenente ci
aveva detto. Poi venne il mio turno col medico che indossava abiti
civili e un camice bianco da laboratorio. Sembrava annoiata. «Devo
togliermi la maglia, signora?» Speravo di potermi tenere i pantaloni
addosso. Era già abbastanza imbarazzante stare in piedi in boxer di
fronte a un medico di sesso maschile, con uno di sesso femminile era
anche peggio. Cioè, non volevo che il mio amico fosse ovviamente felice
di vederla, non so se mi spiego. Ma forse sarebbe stato peggio se non
fosse stato all'altezza della situazione. La dottoressa era carina e io
ero un ragazzo sano, perciò sentii che per rispetto nei suoi confronti
avrei dovuto almeno avere una mezza\ldots{}

«No, soldato, tieni la maglietta. Alza il mento, per favore.» Mi puntò
quella che sembrava una pistola di plastica e cromo sotto l'orecchio
sinistro e premette il grilletto. Non sentii quasi niente. Stessa cosa a
destra. Mise giù la pistola e controllò lo schermo di un computer, ci
tamburellò sopra un paio di volte con le dita e disse: «Pare sia andata
bene. Ti sono state iniettate le nano-macchine kristang che presto
migreranno nel tuo orecchio interno e t'impediranno di provare nausea in
assenza di gravità. Dopo un po' si dissolvono, e non dovresti sentire
alcun effetto collaterale. Se ti sembra che il tuo equilibrio sia
compromesso mentre sei in presenza di gravità, contatta il personale
medico a bordo della nave. Riceverete istruzioni dettagliate una volta
lì».

La porta di fronte a quella da cui ero entrato si aprì e mi fecero cenno
di uscire. Con la gravità non c'erano problemi, era il pensiero che mi
avevano iniettato delle piccole macchine a farmi venire la nausea.
Microscopici robot alieni nuotavano nel mio sangue. Era inquietante.

\hfill\break

Una volta che il nostro gruppo ebbe finito con i medici, le guardie ci
indirizzarono alla nave da trasporto e ci spintonarono senza lasciarci
tempo per un giro turistico. Comunque, non c'era molto da vedere nei
corridoi della stazione: pareti vuote di colore blu chiaro su cui erano
attaccati segnali in diverse lingue umane. La nave kristang era più
interessante, molti pannelli di accesso, corridoi che andavano in ogni
direzione, tubi o condutture lungo i soffitti. Nei muri e nei pavimenti
erano incassate delle maniglie, immagino per quando la nave si fosse
trovata in assenza di gravità. Non si trovava in giro neanche un
kristang, non avevo ancora mai visto una lucertola in persona. Un Marine
prese il comando del nostro gruppo e diresse gli uomini in un
compartimento, le donne in un altro. Nel compartimento in cui finii io
c'erano file e file di letti a castello, impilati a tre a tre. Ogni
gruppo di tre cuccette era attaccato a un braccio in grado di far girare
la pila di novanta gradi quando la nave era sotto propulsione. Il Marine
ci disse che si trattava di una navicella multi-missione kristang,
attualmente configurata per il trasporto di truppe. Eravamo fortunati,
aggiunse, che le cuccette fossero a misura di kristang, che in media
erano un po' più alti degli umani. Alzai la testa e trovai una branda
vuota, ci buttai dentro la mia piccola sacca da viaggio e la fissai con
una cinghia che pensavo fosse lì per quello scopo. La cuccetta non era
male, due metri di lunghezza, a circa un metro e venti centimetri dal
letto di sopra. Era quasi di lusso. Qualcuno era già nel bel mezzo di
una partita a carte: Jesse andò a vedere, mentre io mi lasciai andare di
peso sul letto e selezionai un libro da leggere sul tablet. Stavo
entrando in ansia perché la batteria era al 25\% e sembrava non ci fosse
una presa elettrica da nessuna parte. Avevo l'impressione che fosse
stato stupido portare quella cosa nello spazio con me.

Arrivò un sergente e rimase in piedi sulla porta. Si schiarì la gola per
attirare la nostra attenzione. Mi divertiva vedere che indossava la
stessa giacca mimetica fuori moda che avevo io. «Sono il sergente scelto
Raynor, sono a capo di questo compartimento e degli altri due di fianco
a voi. Tra due ore, questa nave si staccherà e decollerà dall'orbita
terrestre. Dovete essere tutti nelle vostre cuccette con le cinture
allacciate dieci minuti prima della partenza. Se non sapete come usare
le cinghie, lo specialista Edwards ve lo mostrerà. Dopo la partenza
dalla stazione, la gravità artificiale verrà interrotta, quindi
assicuratevi di non lasciare sciolto nulla che possa fluttuare. Questa
nave prenderà velocità per circa un'ora, l'accelerazione è solo al 30\%
della gravità normale, non è come nei razzi della Nasa. Poi l'astronave
farà un salto nell'iperspazio. Potrete uscire dalle vostre cuccette una
volta che i kristang avranno dato il via libera, dopo circa novanta
minuti. Dopo di che, sarete liberi di muovervi; tenete a mente che
saremo a gravità zero. Non voglio che nessun asino gironzoli qua e là a
gravità zero! Una volta che la nave sarà abbastanza lontana dalla Terra,
faremo una serie di salti nell'iperspazio, durante i quali dovrete
starvene nelle vostre cuccette. La nostra destinazione finale è un
pianeta che i kristang hanno adibito a base di addestramento. Se avete
domande, avrete l'opportunità di ottenere risposte più tardi.»

Visto che stavo per rimanere bloccato nella mia cuccetta per più di
un'ora, decisi di andare in bagno. Uscii di nuovo dalla porta e chiesi
alla guardia Marine dove fossero i servizi. M'indicò il fondo del
corridoio, con un sorriso ironico.

I bagni erano interessanti. Non c'erano orinatoi. Si poteva stare seduti
in modo normale, ma erano progettati per l'assenza di gravità, c'erano
cinghie per tenerti in posizione e il sedile ti formava un sigillo sul
fondoschiena. C'era un tubo flessibile con tazze usa e getta per chi
doveva solo fare pipì. Tutte le postazioni avevano targhe di istruzioni
in diverse lingue umane, e non era così complicato come sembrava, il
computer del bagno si prendeva cura di tutto. Segnate un punto per la
tecnologia. A quanto pare, anche i kristang avevano bagni unisex, perciò
facevamo a turno: cinque donne prima, poi cinque uomini. Cazzo, il turno
delle donne durava all'infinito. Ma erano molte meno rispetto agli
uomini, quindi non era un problema. Mentre tornavo al compartimento a me
assegnato, c'era un capitano dell'esercito in piedi nel corridoio che
leggeva qualcosa su un tablet. Feci il saluto e mi fermai a tentare di
parlare con lui: «Mi scusi, signore, come fa a caricare il tablet?».

«In ogni compartimento c'è un set di cuscinetti magnetici, mettici sopra
il tablet e si caricherà. È come un tappetino di ricarica e percepisce
di che tipo di corrente ha bisogno il tuo tablet, in modo da non
fulminarlo. I dispositivi thuranin sono intelligenti.»

«Cos'è un thuranin, signore?»

Il capitano scosse la testa: «I thuranin sono una razza molto evoluta,
guidano la coalizione cui appartengono i kristang. Mi hanno detto che
assomigliano un po' agli alieni dei film di fantascienza: bassi e
ossuti, con grandi teste calve».

Adesso ero in completa confusione: «I kristang non sono al comando?».

«Per quanto riguarda noi umani, sì: per qualsiasi cosa abbia a che fare
con la Terra, i kristang sono al comando. I thuranin sono i
patroni\ldots{} sì, forse è questa la parola giusta, dei kristang.
Comunque, tutto quello che devi sapere è che viaggiare da una stella
all'altra è un processo lento in una nave kristang, quindi si agganciano
alle navi madri thuranin per i grandi salti. E, prima che tu me lo
chieda, no, non ho mai visto un thuranin, siamo a bordo di astronavi
kristang.»

«Sta diventando complicato, signore.»

Lui annuì con simpatia. «Pensala così: i thuranin sono le squadre
maggiori, i kristang sono\ldots{} ah, la squadra satellite in tripla A.
E noi, beh, siamo una specie di squadra semiprofessionistica in questo
momento. Speriamo di essere promossi in A, ma prima dobbiamo dimostrare
quanto valiamo, e abbiamo ancora molta strada da fare. Torna nella tua
cuccetta, soldato, tra poco decolleremo dall'orbita. Riceverai
istruzioni più tardi.»

\hfill\break

Dopo un'ora di leggera accelerazione per allontanarci dalla Terra e
dieci minuti di quella che immaginai fosse la preparazione, la nostra
nave da trasporto saltò. Non riuscivo a vedere niente, legato nella mia
branda. Quando la nave saltò, per una frazione di secondo ebbi la
sensazione che dell'elettricità statica mi stesse strisciando sulla
pelle, poi si tornò a galleggiare in assenza di gravità. Tenni traccia
del tempo sul mio orologio: tra il primo e il secondo salto trascorsero
otto ore e ventuno minuti. Più tardi, appresi che le navi kristang sono
alimentate da reattori a fusione e i loro motori impiegano circa otto
ore per produrre abbastanza energia per un salto. Quel lasso di tempo
tra i salti può rappresentare una grave responsabilità tattica per una
forza in attacco: le navi in attacco, infatti, appena saltate nel
combattimento, possono per lo più saltare via subito se si mettono nei
guai, mentre una nave in difesa, con motori del tutto carichi, può
saltare in una posizione diversa se ha uno sciame di missili in arrivo.
In una battaglia reale, l'intervallo tra i salti non è un problema. Le
navi da guerra mantengono sempre l'energia di riserva per un breve salto
di emergenza. Tuttavia esse non possono saltare in modo sicuro dentro o
fuori dall'iperspazio vicino a un grande campo gravitazionale, come un
pianeta, il campo di salto potrebbe distorcersi e fare a pezzi
l'astronave. Quindi una forza d'attacco deve uscire dall'iperspazio,
lontano da un pianeta, e viaggiare attraverso lo spazio normale per
entrare in orbita. Una volta che le navi attaccanti si avvicinano
abbastanza a un pianeta da usare le armi, si sono spinte così tanto
all'interno del campo di gravità che anche le navi con motori del tutto
carichi non possono saltare in sicurezza. Le navi in difesa, che non
sanno mai dove un attaccante uscirà dall'iperspazio, tendono a volare a
punto fisso vicino ai pianeti, e poi si muovono per intercettare. Suona
come un gioco con troppe regole da memorizzare, ma così è. Per fortuna
per noi umani, non dovevamo pensarci. Non avevamo astronavi e i kristang
non ce ne avrebbero date. Eravamo fanti, combattevamo a terra. I
kristang ci avrebbero portati lì e noi ci saremmo occupati della lotta.
Questo è quello che tutti pensavamo allora, almeno.

Tra un salto e l'altro, l'astronave incrementava per circa due ore, poi
ci davano del tempo libero, quello in cui potevamo uscire dalle nostre
cuccette. Tutti sperimentarono il movimento in assenza di gravità, io
capii dopo qualche minuto che avrebbero fatto meglio a darci dei caschi.
Il sergente Raynor venne a urlarci contro per la nostra stupidità, ma
per lo più lasciava che ci prendessimo la nostra dose di lividi e
ammaccature. L'esperienza è la migliore insegnante.

Uno degli esibizionisti era Jeff Murdock. Era entrato a far parte del
nostro plotone in Nigeria a metà della mia missione lì, quindi non lo
conoscevo bene, ma, ragazzi, affrontava l'assenza di gravità come un
delfino affronta l'acqua. Compiva capriole che mi avrebbero fatto girare
la testa e riusciva a spingersi da un muro e galleggiare attraverso il
compartimento fino a dove voleva andare, nuotando nell'aria. Pendeva da
un letto a testa in giù, usando i piedi per tenersi in posizione, con un
enorme sorriso. «Cazzo, è fantastico! Non vedo l'ora di prendere un po'
di quella biotecnologia kristang. Sarò un super-soldato, amico!»

Noialtri non eravamo così entusiasti di vederlo volare come uno
scoiattolo gigante. «Murdock, non hai bisogno di tecnologia aliena
avanzata per essere un soldato, hai bisogno di un cazzo di miracolo»,
disse Garcia ridendo.

«Sì, amico, quello che ti serve è un cervello, e faresti meglio a
parlarne con il Mago di Oz», fu il commento di Thompson.

«Ehi, vaffanculo!», sparò Murdock di rimando e volò attraverso il
compartimento, girò e atterrò con grazia sul muro opposto, non rimbalzò
nemmeno. «Lo vedrai.»

Ci portarono alla versione kristang di una sala mensa o un refettorio, o
forse si chiamava cambusa, visto che si trovava su una nave?
L'esperienza a gravità zero, fatta nel compartimento dei letti a
castello, fu d'aiuto quando dovemmo trascinarci lungo i corridoi della
nave. La sala mensa aveva la capienza di quasi un centinaio di persone e
c'erano tavoli e sedie, ciò che la rendeva diversa è che quelle sedie
erano saldamente attaccate al pavimento e avevano le cinture di
sicurezza. Il pranzo era un pasto pronto al consumo; gli umani non
potevano mangiare cibo kristang e i kristang non volevano che
incasinassimo le loro cucine a gravità zero. I pasti pronti mi
ricordavano il campeggio di quando ero bambino; non erano male, ma ne
avevo consumati troppi in Nigeria e mangiavo di malavoglia il mio cibo,
scambiando le cose con le persone intorno a me. Cornpone era
super-affamato, ingoiò il suo pasto pronto e mi fregò un pacchetto di
cracker. A metà della pausa pranzo, un maggiore dell'esercito entrò
fluttuando nello scompartimento, volò abilmente fino alla parete opposta
e agganciò i piedi a una cinghia. La faceva sembrare facile. «State
tutti godendovi il cibo degli astronauti?» Strappò una risata alla
folla. «Mi rendo conto che girano un sacco di informazioni sbagliate,
quindi ho intenzione di chiarirvi un po' le idee. Primo, non saremo in
assenza di gravità per tutto il percorso. Quest'astronave kristang starà
agganciata a una nave madre thuranin per la maggior parte del viaggio. I
punti d'attacco sulla nave madre sono hangar con gravità artificiale,
quindi, dopo essere stati attaccati, avremo l'85\% di gravità, che è
normale per i thuranin. Secondo, questo viaggio non durerà quattro anni
e non sarete congelati. Quindi, vi toccherà ascoltare i vostri compagni
di branda russare.» Un'altra risatina dalla folla. «La nave madre
thuranin salterà verso un wormhole, un cunicolo spazio-temporale
artificiale e lo attraverserà tutto fino dall'altra parte, che dista
centinaia di anni luce. La nave madre viaggia tra i wormhole e lascia
che ci portino a destinazione. Saranno diciassette i giorni dalla Terra
alla nostra prima destinazione, un pianeta che chiamiamo Campo Alfa. È
una base di addestramento dove riceverete le vostre nuove armi e
attrezzature e imparerete le regole di ingaggio. I kristang sono pignoli
rispetto alla nostra adesione a esse, quindi rigate dritto mentre siete
lì. Non causeremo problemi ai nostri nuovi alleati. Per quanto riguarda
Campo Alfa, ci sono stato, e tutto quello che dirò ora è che è un posto
eccellente per concentrarci sulla nostra missione di addestramento.»

Questo provocò un gemito d'intesa tra la folla. Tutti capimmo che era un
modo militaresco per parlare di un buco infernale dove non c'era altro
da fare che lavorare e addestrarsi. Guardai Cornpone e alzammo le
spalle. Avevamo visto di peggio e non eravamo nello spazio per una
crociera di piacere. Il maggiore non rispose a nessuna domanda e lasciò
la sala mensa dopo aver finito il suo breve discorso. Mi sentivo molto
meglio sapendo che non saremmo rimasti bloccati a bordo della nave per
quattro anni.

\hfill\break

Dopo che la nave kristang ebbe fatto un quarto salto, dovemmo rimanere
legati nelle nostre cuccette mentre faceva manovra per agganciarsi a una
nave madre thuranin. Ci furono un rumore metallico e una vibrazione. La
gravità tornò lentamente, nell'esultanza di chiunque potessi sentire.
Specie quando annunciarono che per cena avremmo mangiato cibo caldo, ora
che nelle cucine era tornata la gravità. Ci propinarono stufato di
manzo, con molte più patate e verdure che carne, ma era una promessa che
non saremmo sopravvissuti solo con cibi pronti per le successive tre
settimane. I biscotti erano caldi e appena sfornati. Il buon cibo è
importante per il buon morale, l'esercito di sicuro lo sapeva.

Mi piacciono i biscotti caldi appena sfornati.

Nel caso ve lo steste chiedendo.

Il cargo stellare thuranin, la nave madre, o comunque la si chiami, fece
una serie di salti di cui persi il conto, perché alcuni avvennero mentre
dormivo. Certi a bordo stavano cercando di tenere un registro amatoriale
di tutte le manovre: salti, tempo tra i salti, tempo e forza di
accelerazione. Per quanto mi riguarda, persi il conto dopo il secondo
salto e pensai che ci dovessero essere persone qualificate che tenevano
il conto per noi. Gli altoparlanti ci svegliarono intorno alle tre di
notte, quando la nave si stava avvicinando al wormhole, e avrei voluto
avere un più ampio preavviso, perché la mia vescica mi stava già
mandando segnali di avvertimento. L'idea di cadere in un wormhole, coi
miei atomi fatti a pezzi e rimontati, o qualsiasi cosa accadesse, mi
spaventava. Ero in posizione rigida nella mia cuccetta, in attesa
dell'ignoto, quando l'altoparlante annunciò: «Transizione wormhole
completata. Allacciate le cinture per una normale manovra spaziale».
Cazzo. Non solo non avrei avuto bisogno di preoccuparmi del passaggio
nel wormhole, ma non mi ero neanche accorto di quando era avvenuto. Ci
furono altri salti, altra attesa fra un salto e l'altro e un sacco di
noia. Ci allenavamo a turno, utilizzando attrezzature portate dalla
Terra, tra cui tapis roulant e cyclette. Mi piaceva correre, ma odiavo
il tapis roulant, era così noioso. Anche competere con il tizio sul
tapis roulant accanto a me non era d'aiuto, dopo i primi due giorni.
Giocavamo sui tablet, ci esercitavamo, leggevamo libri e guardavamo
film. Essere a bordo della nave era un po' come doveva essere per i
soldati della seconda guerra mondiale attraversare l'Atlantico e il
Pacifico su lente navi da trasporto. Non fosse che, a bordo di una nave
da trasporto della seconda guerra mondiale, quei ragazzi avevano la
possibilità di salire sul ponte, prendere aria fresca e vedere il sole e
l'orizzonte.

Avevamo cibo decente e abbastanza spazio per muoverci in giro per la
nave, a parte le aree off limits per gli umani. Quello che non avevamo
era qualcosa da guardare fuori dal finestrino, se mai ci fossero stati
finestrini. E non avevamo notizie dalla Terra, né contatti con essa. In
Nigeria potevo accedere a internet, inviare e-mail e video-chattare con
i miei genitori e amici, almeno la maggior parte dei giorni. Stavolta
non avremmo avuto contatti fino al ritorno, molto, molto lontano. Una
cosa che imparammo è che nemmeno i thuranin, per quanto superiori dal
punto di vista tecnologico, avevano mezzi di comunicazione più veloci
della luce, a eccezione dei droni che trasportavano messaggi saltando.
Un po' come un piccione viaggiatore ad alta tecnologia.

Dopo un po', mi sentivo come se fossimo turisti a basso budget bloccati
a bordo di una nave da crociera di merda, non sembrava di essere
nell'esercito. Dopo qualche altro salto nell'iperspazio, di cui non
cercai nemmeno di tenere il conto, la nave madre sbucò fuori da qualche
sistema stellare, o almeno così ci dissero. La nave kristang si ritrovò
a gravità zero, si staccò e saltò quattro volte. Infine, facemmo manovre
in orbita sopra la nostra destinazione. Tutto sommato, il mio primo
viaggio interstellare consistette per lo più nel restare bloccato in un
compartimento senza finestrini. Non lo consiglierei per una vacanza, a
meno che non siate disperati.

Gettammo un'occhiata al pianeta sottostante, non era invitante. La Terra
dall'orbita appare blu, verde e marrone chiaro dove ci sono i deserti e
le praterie, bianca ai poli dove c'è il ghiaccio. Questo posto appariva
marrone e rosso, e anche i luoghi dove scorreva l'acqua non erano di un
bel blu. L'unico polo che vidi presentava uno strato anemico di bianco
sporco, come la neve accanto alla strada in aprile, quando è sporca e
crostosa e tutti vorrebbero che si sciogliesse e se ne andasse, cazzo.

A gruppi, ci ammassammo nella sala mensa per avere informazioni. Il
tenente Gonzalez stava in piedi su una sedia, e aspettò fino a quando
tutti si furono sistemati nel compartimento. «Benvenuti su Campo Alfa! I
kristang non volevano dirci il nome che usavano per questo pianeta,
quindi per noi era solo ``Alfa''.»

Un tizio in prima fila chiese: «Ma sappiamo dove siamo, giusto, tenente?
Voglio dire, qualcuno l'ha capito?».

«Sì», sorrise Gonzalez, «la squadra d'avanscoperta ha scattato foto del
cielo notturno e gli astronomi sulla Terra hanno calcolato il punto dal
quale doveva essere stata fatta l'inquadratura. Penso che i kristang si
stessero divertendo con noi tenendolo segreto, ma dovevano sapere che
saremmo stati abbastanza svegli da scoprire la posizione di questo
pianeta. Siamo ancora nel Braccio di Orione della galassia, la Via
Lattea, a 1.223 anni luce dalla Terra. È abbastanza lontano e perciò non
si può vedere il nostro Sole da qui senza un telescopio. Mi hanno detto
che è una stella insignificante comunque.»

«Avete scrutato tutti dal finestrino il luogo in cui ci addestreremo?
Campo Alfa non è il tipo di posto dove andreste in vacanza. È il tipo di
posto dove i soldati sono addestrati e formati in un esercito che i
ruhar impareranno a temere. I kristang hanno destinato questo pianeta
come base di addestramento e preparazione per gli esseri umani per tre
ragioni. Prima, nessun altro vuole questo posto: è asciutto, mi hanno
detto che la temperatura può raggiungere i centodieci gradi durante il
giorno e scendere fino a venti gradi sottozero la notte. Questa è la
parte bella del pianeta, dove si trova la base di addestramento. La luce
stellare è di un bizzarro colore blu-bianco e vi sbuccerà la pelle dalle
ossa se non siete protetti. Seconda, poiché questo posto è così orrendo,
è un ottimo posto per allenarsi. Non tutti i pianeti in cui andremo
saranno come la Terra, qui imparerete a usare l'equipaggiamento fornito
dai kristang e a sopravvivere in condizioni difficili. Se non siete in
grado di combattere in modo efficiente qui, non abbiamo bisogno di voi.
E terza, i kristang hanno scelto questo posto perché pensano che i ruhar
non lo conoscano. Eppure\ldots»

«Signore», chiese qualcuno, «perché il pianeta appare tutto marrone e
bruciacchiato?»

«La stella qui è una variabile lenta, singhiozza di tanto in tanto e
aumenta la sua potenza. Quando questo accade, brucia l'intero sistema,
poi si calma di nuovo.»

Singhiozzo? Alla faccia del singhiozzo. A disagio, ci spostammo in massa
e tra di noi si diffuse un mormorio sommesso.

«Non preoccupatevi di questo, i kristang conoscono il ciclo della
stella, non si riaccenderà per altri dieci o cinquantamila anni»,
aggiunse in fretta Gonzalez. Dai dieci ai cinquantamila anni: sembrava
che i kristang non avessero la più pallida idea di quando sarebbe
esplosa la stella. «I vostri pacchetti contengono tutte le informazioni
che vi servono su questo pianeta, quindi citerò solo i punti salienti.
C'è vita laggiù, l'animale più grande è un insetto lungo circa dieci
centimetri. Ha pinze affilate ed è velenoso, ma il veleno non è fatale
per la biochimica umana, provoca solo una lieve eruzione cutanea, a
volte, come reazione allergica. La maggior parte della gente non
reagisce affatto. Le uniche cose pericolose per noi laggiù sono il
clima, che può diventare molto caldo e molto freddo, e la nostra
stupidità. Usate il buon senso, ricordate che siete qui per
l'addestramento e leggete i regolamenti: sono scritti per tenervi in
vita e pronti a combattere. Ci raggrupperemo in unità quando saremo
sulla superficie, fino ad allora, seguite le istruzioni.»